\documentclass{article}

\usepackage{amsmath,amssymb}
\usepackage{xurl}
\usepackage[hidelinks]{hyperref}
\setlength{\emergencystretch}{2em}

% Shared commands for notes in docs/paper-gaps/.

\DeclareUnicodeCharacter{2080}{\ifmmode{}_{0}\else\textsubscript{0}\fi}
\DeclareUnicodeCharacter{2081}{\ifmmode{}_{1}\else\textsubscript{1}\fi}
\DeclareUnicodeCharacter{2082}{\ifmmode{}_{2}\else\textsubscript{2}\fi}
\DeclareUnicodeCharacter{2099}{\ifmmode{}_{n}\else\textsubscript{n}\fi}

\newcommand{\Lean}{\textsc{Lean}}
\ExplSyntaxOn
\NewDocumentCommand{\leanid}{m}{
  \ifmmode
    \textsf{\detokenize{#1}}
  \else
    \group_begin:
    \urlstyle{sf}
    \tl_set:Nn \l_tmpa_tl {#1}
    \tl_replace_all:Nnn \l_tmpa_tl {\_} {_}
    \exp_args:NV \path \l_tmpa_tl
    \group_end:
  \fi
}
\ExplSyntaxOff
\providecommand{\tr}{\operatorname{tr}}
\newcommand{\tnleanIssueBase}{https://github.com/LionSR/TNLean/issues/}
\newcommand{\tnleanPrBase}{https://github.com/LionSR/TNLean/pull/}
\newcommand{\tnleanDocBase}{https://sirui-lu.com/TNLean/docs/}
\newcommand{\ghissue}[1]{\href{\tnleanIssueBase#1}{GitHub issue \##1}}
\newcommand{\ghpr}[1]{\href{\tnleanPrBase#1}{PR \##1}}
\newcommand{\doclink}[1]{\href{\tnleanDocBase#1.html}{\path{#1}}}

% Dirac notation and the identity operator, as in blueprint/src/macros/common.tex.
% \providecommand keeps note-local definitions authoritative.
\usepackage{dsfont}
\providecommand{\ket}[1]{|#1\rangle}
\providecommand{\bra}[1]{\langle #1|}
\providecommand{\braket}[2]{\langle #1 | #2 \rangle}
\providecommand{\ketbra}[2]{|#1\rangle\!\langle #2|}
\providecommand{\Id}{\mathds{1}}
\providecommand{\one}{\mathds{1}}
\providecommand{\C}{\mathbb{C}}
\providecommand{\MN}[1]{M_{#1}(\C)}
\providecommand{\MD}{M_D(\C)}

% External referencing for paper sources.  The build helper writes sanitized
% auxiliary files under build/paper-aux/ and excludes duplicated source labels.
\usepackage{xr}

% Import a generated paper auxiliary file with a caller-supplied label prefix.
% The candidate paths support builds from docs/paper-gaps/, the repository root,
% and build/paper-aux/ itself.
\newcommand{\importpaperaux}[2]{%
  \IfFileExists{../../build/paper-aux/#2.aux}{%
    \externaldocument[#1]{../../build/paper-aux/#2}%
  }{%
    \IfFileExists{build/paper-aux/#2.aux}{%
      \externaldocument[#1]{build/paper-aux/#2}%
    }{%
      \IfFileExists{#2.aux}{%
        \externaldocument[#1]{#2}%
      }{}%
    }%
  }%
}

% General external reference macros
\newcommand{\paperref}[2]{\ref{#1-#2}}
\newcommand{\papereqref}[2]{(\ref{#1-#2})}

% CPSV16 source (1606.00608 / MPDO-22-12-17-2.tex) external document and shortcuts
\importpaperaux{cpsv16-}{1606.00608/MPDO-22-12-17-2-tnlean}
\newcommand{\cpsvref}[1]{\paperref{cpsv16}{#1}}
\newcommand{\cpsveqref}[1]{\papereqref{cpsv16}{#1}}

% Verdict marker: \gapnote{<kind>}{<status>}. Kind is the policy
% classification (clarification, local-correction, scope-restriction,
% unfaithful, false-source, open-gap); status is open, wip (elimination
% underway), resolved, or historical. Machine-read by the site index;
% prints a verdict line.
\newcommand{\gapnote}[2]{\par\noindent\textbf{Verdict:} #1 (#2).\par}

\importpaperaux{fbc25-}{2502.20257/main-tnlean}

\title{Source Boundary for the MPU Gauss Laws and Their Commuting Projectors}
\author{}
\date{2026-08-29}

\begin{document}
\maketitle
\gapnote{open-gap}{open}

\section{Scope}

Proposition~\paperref{fbc25}{prop:gausslaws} of
Ref.~\cite[lines~3378--3499]{FrancoRubio2025MPUGauging} defines a local
unitary representation $\mathcal G$ and the corresponding Gauss-law
projectors $P_j$.  Proposition~\paperref{fbc25}{prop:comm_proj} then states that the
projectors commute if and only if the MPU symmetry is nonanomalous, and refers
to Section~3.2 of Ref.~\cite{Seifnashri2024LSM} for the adjacent-projector
calculation \cite[lines~3602--3655]{FrancoRubio2025MPUGauging}.

This note fixes the mathematical meaning of that citation before either
proposition is formalized.  It records the oriented local Hilbert space, the
multiplication convention, the hypotheses used by the three-defect
calculation, and the exact boundary of the existing formal library.  It does
not replace the cited tensor calculation by an abstract cocycle argument.
The Seifnashri line ranges below refer to
\emph{Lieb-Schultz-Mattis anomalies as obstructions to gauging
(non-on-site) symmetries},
\href{https://arxiv.org/abs/2308.05151v3}{arXiv:2308.05151v3}, source file
\path{LSMcocycle.tex}.  This version is distributed under the arXiv
non-exclusive distribution license 1.0.
The formalization prerequisite is tracked by \ghissue{7339}.

\section{Oriented local Hilbert space}

Let $G$ be a finite group and let the matter space at site $j$ be
$\mathcal H_j\cong\mathcal H$.  In the Hamiltonian construction of
Ref.~\cite[Section~3.2]{Seifnashri2024LSM}, the chain is periodic, a defect
$g_j$ lies on the oriented link $(j,j+1)$, and the gauge basis on that link is
written $\ket{g_j}$.  The local transformation at $j$ uses two matter factors
and the two links adjacent to $j$.  After the harmless reordering of tensor
factors used in Ref.~\cite[lines~3380--3386]{FrancoRubio2025MPUGauging}, its
space is
\begin{align}
    \mathcal K_j
        &=\mathcal H_j\otimes\mathcal H_{j+1}
          \otimes\mathbb C[G]_{(j-1,j)}
          \otimes\mathbb C[G]_{(j,j+1)}.
    \label{eq:fbc25_gauss_oriented_local_space}
\end{align}
The first gauge label is therefore incoming and the second is outgoing.  If
their values are $a$ and $b$, a transformation by $g$ sends
\begin{align}
    (a,b)&\longmapsto(ag^{-1},gb).
    \label{eq:fbc25_gauss_oriented_link_action}
\end{align}
Thus the incoming link carries right multiplication by $g^{-1}$, while the
outgoing link carries left multiplication by $g$.  Reversing these roles, or
replacing either action by multiplication on the other side, changes the
representation law.

Seifnashri first uses defect Hilbert spaces which may depend on all link
labels.  For a discrete group, the formula used here is obtained after the
explicit simplifying assumption that every defect Hilbert space is an
identical copy of the untwisted matter space.  The extended space then becomes
\begin{align}
    \mathcal H_{\mathrm{ext}}
        &=\left(\bigotimes_{j=1}^{L}\mathcal H_j\right)
          \otimes
          \left(\bigotimes_{j=1}^{L}\mathbb C[G]_{(j,j+1)}\right),
    \label{eq:fbc25_gauss_extended_space}
\end{align}
with one matter factor per site and one gauge factor per oriented link
\cite[Section~3.2, arXiv source lines~1637--1678]{Seifnashri2024LSM}.

\section{The local representation and its multiplication convention}

For $a,b,c,d\in G$ with $ab=cd$, let
$\lambda_{a,b}^{c,d}$ be the unitary comparison of the two defect-pair
factorizations.  The MPU paper defines it from right fusion operators by
$\lambda_{a,b}^{c,d}=(\lambda^R_{c,d})^\dagger\lambda^R_{a,b}$ and proves the
composition law
\begin{align}
    \lambda_{a,b}^{r,s}
        &=\lambda_{c,d}^{r,s}\lambda_{a,b}^{c,d},
        &ab&=cd=rs.
    \label{eq:fbc25_gauss_lambda_composition}
\end{align}
These are the comparison operators of
Ref.~\cite[lines~2782--2821]{FrancoRubio2025MPUGauging}; they are not new
fusion tensors.

On the oriented space in
Eq.~\ref{eq:fbc25_gauss_oriented_local_space}, the local operator is
\begin{align}
    \mathcal G^j_g
        &=\sum_{a,b\in G}
          \lambda_{a,b}^{ag^{-1},gb}
          \otimes\ketbra{ag^{-1},gb}{a,b}.
    \label{eq:fbc25_gauss_local_representation}
\end{align}
The superscript $j$ records the support and is suppressed in the displayed
formula \papereqref{fbc25}{eq:mcG} of
Ref.~\cite[lines~3380--3386]{FrancoRubio2025MPUGauging}.  Applying first $h$
and then $g$ to one gauge sector gives
\begin{align}
    (a,b)
        &\overset{\mathcal G^j_h}{\longmapsto}(ah^{-1},hb)
         \overset{\mathcal G^j_g}{\longmapsto}
         (ah^{-1}g^{-1},ghb)
         =(a(gh)^{-1},(gh)b),
    \label{eq:fbc25_gauss_sector_multiplication}\\
    \lambda_{ah^{-1},hb}^{ah^{-1}g^{-1},ghb}
      \lambda_{a,b}^{ah^{-1},hb}
        &=\lambda_{a,b}^{a(gh)^{-1},(gh)b}.
    \label{eq:fbc25_gauss_material_multiplication}
\end{align}
Equation~\ref{eq:fbc25_gauss_material_multiplication} is precisely
Eq.~\ref{eq:fbc25_gauss_lambda_composition}.  Hence the convention is
\begin{align}
    \mathcal G^j_g\mathcal G^j_h
        &=\mathcal G^j_{gh},
    &
    (\mathcal G^j_g)^\dagger
        &=\mathcal G^j_{g^{-1}}.
    \label{eq:fbc25_gauss_representation_law}
\end{align}
This agrees with the explicit convention in
Ref.~\cite[arXiv source lines~1691--1734]{Seifnashri2024LSM}.  Unitarity is
not a consequence of the label permutation alone: every nonzero block
$\lambda_{a,b}^{ag^{-1},gb}$ must be unitary.  In Seifnashri's construction
this comes from unitary defect fusion operators $\lambda^j(g,h)$; the
link-extended operators $\widetilde\lambda^j(g,h)$ are partial operators and
are explicitly not unitary \cite[arXiv source lines~709--744 and
1680--1729]{Seifnashri2024LSM}.

The Gauss-law projector is the group average
\begin{align}
    P_j
        &=\frac{1}{|G|}\sum_{g\in G}\mathcal G^j_g.
    \label{eq:fbc25_gauss_projector}
\end{align}
Equation~\ref{eq:fbc25_gauss_representation_law} makes
$P_j^2=P_j=P_j^\dagger$.  This is the local projection onto the invariant
subspace; the additional constraint is $P_j=1$ on physical states.

\section{The three-defect associator calculation}

The commuting-projector proof uses more than the scalar $3$-cocycle equation.
Let $\omega\colon G^3\to\mathbb C^\times$ be the anomaly $3$-cocycle defined by
the MPU fusion tensors in
Ref.~\cite[lines~1506--1545]{FrancoRubio2025MPUGauging}.
Let $T_{g(hk)}$ and $T_{(gh)k}$ denote the two unitary fusion paths for three
consecutive defects in Equation~\papereqref{fbc25}{eq:assoc} of
Ref.~\cite[lines~3602--3655]{FrancoRubio2025MPUGauging}.  In the notation of
its Appendix~\paperref{fbc25}{app:associativity}
\cite[lines~6040--6961]{FrancoRubio2025MPUGauging}, they satisfy
\begin{align}
    T_{g(hk)}
        &=F(g,h,k)T_{(gh)k}.
    \label{eq:fbc25_gauss_three_defect_associator}
\end{align}
The left path contains the movement operator
$w_R^g=\lambda^R_{g,e}$ between the fusions of $(h,k)$ and $(g,hk)$; the
right path fuses $(g,h)$ and then $(gh,k)$.  The appendix computes the two
tensor networks and obtains
\begin{align}
    F(g,h,k)&=\omega(g,h,k).
    \label{eq:fbc25_gauss_associator_is_omega}
\end{align}
This is a tensor-level result, proved at lines~6040--6962 of
Ref.~\cite{FrancoRubio2025MPUGauging}; it must not be inferred merely from
the existence of an independently supplied scalar cochain.

To compare $P_jP_{j-1}$ with $P_{j-1}P_j$, the cited calculation expands both
products in defect sectors.  The bottom labels
$g_1,g_2,g_3$ and top labels $h_1,h_2,h_3$ have equal total product.  A
unitary fusion of the three top defects converts equality of the two expanded
operators into equality of the two three-defect networks.  Unitarity then
cancels every fusion bubble through
$\lambda^j(g,h)(\lambda^j(g,h))^\dagger=1$.  Applying an $F$-move on each
side, followed by two further $F$-moves on the right, gives exactly
\begin{align}
    \frac{F(h_1,h_2h_3g_3^{-1},g_3)}
         {F(g_1,g_2,g_3)}
        &=
      \frac{F(h_1,h_2,h_3)}
           {F(g_1,g_1^{-1}h_1h_2,h_3)}.
    \label{eq:fbc25_gauss_adjacent_commutation_condition}
\end{align}
This is the calculation at arXiv source lines~1745--1955 of
Ref.~\cite{Seifnashri2024LSM}, relabelled as in lines~3607--3650 of
Ref.~\cite{FrancoRubio2025MPUGauging}.  The denominators are legitimate
because the fusion paths are unitary, so $F$ has unit modulus in the chosen
gauge.

If the projectors commute, set
$h_1=e$, $h_2=g_1$, and $h_3=g_2g_3$ in
Eq.~\ref{eq:fbc25_gauss_adjacent_commutation_condition}.  One obtains
\begin{align}
    F(g_1,g_2,g_3)
        &=\frac{F(g_1,e,g_2g_3)F(e,g_1g_2,g_3)}
        {F(e,g_1,g_2g_3)}.
    \label{eq:fbc25_gauss_specialized_associator}
\end{align}
The standing identity gauges give
$F(e,g,h)=F(g,e,h)=1$, hence
$F(g_1,g_2,g_3)=1$.  Conversely, pointwise $F=1$ makes
Eq.~\ref{eq:fbc25_gauss_adjacent_commutation_condition} automatic.  This is
the actual ``if and only if'' calculation.  The scalar cocycle identity does
not by itself imply either direction.

\section{Gauge-class triviality versus a pointwise trivial representative}

The MPU paper calls a representation nonanomalous when the cohomology class
of $\omega$ is trivial.  This is an assertion about the orbit under fusion
gauges, not the pointwise equation $\omega(g,h,k)=1$ for an arbitrary
representative.  Proposition~\paperref{fbc25}{prop:unitarity_and_gauges} of
Ref.~\cite{FrancoRubio2025MPUGauging} first constructs a compatible
fusion/action gauge in which all fusion operators are unitary and, in the
nonanomalous case, the selected representative satisfies
\begin{align}
    \omega(g,h,k)&=1
        &&\text{for every }g,h,k\in G.
    \label{eq:fbc25_gauss_pointwise_trivial_gauge}
\end{align}
Only after this choice may
Eq.~\ref{eq:fbc25_gauss_associator_is_omega} and
Eq.~\ref{eq:fbc25_gauss_specialized_associator} identify nonanomalous with
commuting projectors.  Seifnashri makes the same distinction: the direct
calculation says $F^j=1$, while the obstruction is the phase-redefinition
class $[F^j]$ \cite[arXiv source lines~1541--1619 and
1956--1960]{Seifnashri2024LSM}.

The existing scalar API distinguishes class-triviality from the three
identity-axis normalization conditions, but only algebraically over
$\mathbb C^\times$.  Both scalar $2$-cochains and scalar $3$-cochains take
values in $\mathbb C^\times$ (the type \leanid{Units}~$\mathbb C$);
\leanid{TNLean.Algebra.ScalarThreeCochain.IsTrivialGaugeClass} expresses
class-triviality, while
\leanid{TNLean.Algebra.ScalarThreeCochain.IsNormalized} expresses the three
identity-axis normalizations.  No present predicate restricts these units to
the unit-modulus subgroup $U(1)$.  Consequently, this API does not yet encode
faithfully the phase-valued fusion gauge required by
Proposition~\paperref{fbc25}{prop:unitarity_and_gauges} of
Ref.~\cite[lines~3268--3346]{FrancoRubio2025MPUGauging}; that phase/unitarity
restriction remains an explicit formal boundary.  The API also does not
package an $H^3$ quotient or attach a scalar cochain to fusion tensors.
\footnote{Formal declarations are in
\doclink{TNLean/Algebra/ScalarThreeCocycle}.  The corrected three-axis
normalization is documented in
\doclink{docs/paper-gaps/fbc25_three_cocycle_normalization_typo}.}

\section{Finite-chain support convention}

Seifnashri's Section~3.2 uses a periodic chain of $L$ sites.  Earlier, the
defect width is defined independently of the system size, and the discussion
then fixes width $k=1$ so that adjacent defects and fusion operators are
unambiguous \cite[arXiv source lines~555--584 and 602--760]
{Seifnashri2024LSM}.  With this convention, $\mathcal G^j_g$ is supported on
matter sites $j,j+1$ and gauge links $(j-1,j),(j,j+1)$.  Thus projectors whose
supports are separated commute by tensor-factor locality, and the only
load-bearing overlap is the adjacent pair $P_j,P_{j-1}$.

The adjacent calculation uses three distinct matter sites
$j-1,j,j+1$ and three distinct gauge links from $(j-2,j-1)$ through
$(j,j+1)$.  Thus $L\geq3$ is the natural stable finite-periodic convention.
For $L=3$, the two endpoint names $j-2$ and $j+1$ coincide, but the three
matter factors and three oriented link factors used by the operators remain
distinct; every pair of different projectors is adjacent.  For $L=2$, the
three-link sector expansion identifies gauge factors and is not the cited
calculation.  The papers do not separately verify this wrapped case.
Consequently, a theorem restricted to $L\geq3$ does not by itself formalize an
unqualified all-length reading that includes $L=2$; that case requires a
separate calculation.

The current finite-region API supplies identity-tensor inclusions and
commutation for disjoint subsets of $\mathbb Z$.
\footnote{Formal declarations:
\leanid{SpinChain.LocalAlgebra},
\leanid{SpinChain.localInclusion}, and
\leanid{SpinChain.localInclusion\_commute\_of\_disjoint} in
\doclink{TNLean/QCA/LocalAlgebra} and
\doclink{TNLean/QCA/DisjointSupport}.}
It does not yet represent the heterogeneous oriented factors in
Eq.~\ref{eq:fbc25_gauss_oriented_local_space}, nor cyclic wrap-around.
Conversely, \leanid{MPOTensor.IsMPU} and
\leanid{MPOTensor.IsMPU.mpo\_mem\_unitaryGroup} describe homogeneous periodic
MPO operators for every $N>1$, but not an arbitrary-boundary local operator
with separate matter and link factors.
\footnote{Formal declarations are in
\doclink{TNLean/MPS/MPU/Basic}.}
Neither surface is by itself a definition of $\mathcal G^j_g$ or $P_j$.

\section{Existing ownership boundaries}

The remaining work should use the following existing declarations and issue
boundaries.

\begin{enumerate}
    \item The comparison operators $\lambda_{a,b}^{c,d}$ and
    Eq.~\ref{eq:fbc25_gauss_lambda_composition} belong to \ghissue{7334}.
    That issue uses the paper's definition
    $(\lambda^R_{c,d})^\dagger\lambda^R_{a,b}$ and standard matrix-unitary
    multiplication.

    \item Unitary $\lambda^L,\lambda^R$ and the gauge in
    Eq.~\ref{eq:fbc25_gauss_pointwise_trivial_gauge} belong to
    \ghissue{7336}.  That issue owns the trace normalization and the
    compatible fusion/action gauge.

    \item The fusion operators and anomaly data needed for the associator
    comparison belong to \ghissue{7322}.  The tensor theorem $F=\omega$ in
    Eq.~\ref{eq:fbc25_gauss_associator_is_omega} belongs to \ghissue{7359},
    which depends on \ghissue{7322}.  The scalar must be derived from the two
    MPU fusion trees.

    \item The identity and inverse gauges used in
    Eq.~\ref{eq:fbc25_gauss_specialized_associator} belong to
    \ghissue{7325}.  Their normalized identity cases should be reused rather
    than postulated again.

    \item Scalar class and normalization are already expressed by
    \leanid{TNLean.Algebra.ScalarThreeCochain.CohomologousTo},
    \leanid{TNLean.Algebra.ScalarThreeCochain.IsTrivialGaugeClass}, and
    \leanid{TNLean.Algebra.ScalarThreeCochain.IsNormalized}.

    \item Disjoint finite support is already expressed by
    \leanid{SpinChain.localInclusion\_commute\_of\_disjoint}.  It closes only
    the separated-support part after a faithful oriented embedding has been
    given.
\end{enumerate}

Matrix unitarity itself is already expressed by
\leanid{Matrix.unitaryGroup}, with the two equations exposed by
\leanid{Matrix.mem\_unitaryGroup\_iff} and
\leanid{Matrix.mem\_unitaryGroup\_iff'}.  QICLean supplies the matrix
$C^*$-algebra infrastructure used by the finite-region API, but a repository
search finds no defect-fusion, Gauss-law, or oriented matter--link
representation that would make
Eq.~\ref{eq:fbc25_gauss_local_representation} redundant.

\section{Conclusion}

The citation fixes a left representation
$\mathcal G^j_g\mathcal G^j_h=\mathcal G^j_{gh}$ on two matter factors and
the incoming/outgoing gauge links, with the incoming label transformed by
right multiplication by $g^{-1}$ and the outgoing label by left multiplication
by $g$.  The commuting-projector result depends on the full adjacent
three-defect tensor calculation, unitary bubble cancellation, the theorem
$F=\omega$, the normalized identity gauges, and a subsequent choice of a
pointwise trivial representative of a trivial anomaly class.

The open formal boundary is therefore not scalar cohomology.  It is the
construction and oriented finite-chain placement of the unitary comparison
operators, followed by the exact calculation in
Eq.~\ref{eq:fbc25_gauss_adjacent_commutation_condition}.  The stable-range
support argument is $L\geq3$; an unrestricted finite-chain theorem also owes
the $L=2$ wrapped case.  No implementation of
Proposition~\paperref{fbc25}{prop:comm_proj} of
Ref.~\cite[lines~3602--3655]{FrancoRubio2025MPUGauging} is claimed here.

\bibliographystyle{alpha}
\bibliography{references}

\end{document}
